\chapter{יישומים}

\section{\textenglish{BERT} ומידול שפה דו-כיווני}

מודל ה-\textenglish{BERT} (\textenglish{Bidirectional Encoder Representations from Transformers}) \cite{devlin2019bert}
מהווה נקודת מפנה מכרעת בתחום עיבוד השפה הטבעית. בשונה ממודלים קודמים שעיבדו את הרצף בכיוון יחיד --- משמאל לימין או מימין לשמאל --- מאפשרת ארכיטקטורת ה-\textenglish{Transformer} הדו-כיווני לכל אסימון להתייחס לכל שאר האסימונים ברצף בו-זמנית, הן בהקשרו השמאלי והן בהקשרו הימני.

האימון המקדים של \textenglish{BERT} מבוסס על שתי משימות עיקריות. המשימה הראשונה היא מידול שפה מוסווה (\textenglish{Masked Language Modelling}, \textenglish{MLM}): בשלב זה, כ-15\% מהאסימונים ברצף הקלט מוחלפים באסימון מיוחד \textenglish{[MASK]}, והמודל מתאמן לחזות את האסימון המקורי על סמך ההקשר הדו-כיווני המלא. פונקציית ההפסד של משימה זו מוגדרת כ:

\begin{equation}
\mathcal{L}_{\text{MLM}} = -\sum_{i \in \mathcal{M}} \log P\!\left(x_i \mid \tilde{\mathbf{x}}; \theta\right)
\end{equation}

כאשר \(\mathcal{M}\) הוא קבוצת מיקומי האסימונים המוסווים, \(\tilde{\mathbf{x}}\) הוא הרצף לאחר הסוואה, ו-\(\theta\) הם פרמטרי המודל. המשימה השנייה היא ניבוי המשפט הבא (\textenglish{Next Sentence Prediction}, \textenglish{NSP}): המודל מקבל זוג משפטים ומסווג האם המשפט השני אכן עוקב אחר הראשון בטקסט המקורי. משימה זו נועדה להקנות למודל יכולת הבנה של קשרים בין-משפטיים, דבר שחיוני למשימות כגון מענה לשאלות (\textenglish{Question Answering}) וסיק מסקנות טבעי (\textenglish{Natural Language Inference}).

ארכיטקטורת \textenglish{BERT} מבוססת על ערימת מקודדים (\textenglish{encoder stack}) בלבד. הגרסה הבסיסית, \textenglish{BERT-Base}, כוללת 12 שכבות \textenglish{Transformer}, 12 ראשי קשב, וממד הסתר \(d_{\text{model}} = 768\), המסתכם בכ-110 מיליון פרמטרים. גרסת \textenglish{BERT-Large} מכילה 24 שכבות וממד הסתר \(d_{\text{model}} = 1024\), המסתכם בכ-340 מיליון פרמטרים. האימון בוצע על קורפוסים גדולים הכוללים את \textenglish{BooksCorpus} ואת \textenglish{Wikipedia} האנגלית. מחקרים מאוחרים יותר הראו כי משימת ה-\textenglish{NSP} תורמת פחות ממה שהניחו תחילה, ומודלים כגון \textenglish{RoBERTa} השמיטוה ושיפרו את הביצועים בהתאם. גם \textenglish{ELECTRA} \cite{clark2020electra} הציע חלופה יעילה יותר לאימון מקדים, בה המודל מאומן להבחין בין אסימונים אמיתיים לאסימונים שנוצרו על ידי רשת גנרטיבית קטנה.

תרומתו המרכזית של \textenglish{BERT} היא ייצוג הקשרי עמוק (\textenglish{deep contextual representation}): וקטור ההטמעה של כל אסימון תלוי בכל שאר הרצף, כך שמילה פולימסמית כגון ``בנק'' תקבל ייצוג שונה בהקשר פיננסי לעומת הקשר גיאוגרפי. תכונה זו הפכה את \textenglish{BERT} לבסיס לכיוונון עדין במגוון רחב של משימות.

\section{\textenglish{GPT} ומודלים גנרטיביים}

בניגוד לגישת ה-\textenglish{BERT} המבוססת על מקודד דו-כיווני, משפחת מודלי ה-\textenglish{GPT} (\textenglish{Generative Pre-trained Transformer}) \cite{radford2019language} אימצה ארכיטקטורת מפענח (\textenglish{decoder-only}) המאומנת על ידי מידול שפה סיבתי (\textenglish{Causal Language Modelling}, \textenglish{CLM}). פונקציית ההפסד מוגדרת כ:

\begin{equation}
\mathcal{L}_{\text{CLM}} = -\sum_{t=1}^{T} \log P\!\left(x_t \mid x_1, x_2, \ldots, x_{t-1}; \theta\right)
\end{equation}

ארכיטקטורת המפענח מממשת מסכת קשב עתידית (\textenglish{causal attention mask}) המונעת מכל אסימון לצפות באסימונים הבאים אחריו, ובכך שומרת על תכונת הסיבתיות הנדרשת לאימון גנרטיבי. \textenglish{GPT-2} \cite{radford2019language} הדגים לראשונה כי מודל שפה גנרטיבי גדול מסוגל לבצע משימות שונות ללא כל כיוונון עדין. המודל \textenglish{GPT-3} \cite{brown2020language} הרחיב גישה זו ל-175 מיליארד פרמטרים, והדגים יכולות \textenglish{in-context learning} מרשימות. \textenglish{LLaMA} \cite{touvron2023llama} הוכיח כי ניתן להשיג ביצועים מתחרים באמצעות אימון ממושך על קורפוסים ציבוריים ושיפורים ארכיטקטורליים כגון \textenglish{RoPE} ו-\textenglish{SwiGLU}.

\section{כיוונון עדין והעברת למידה}

האימון המקדים מניב ייצוגים כלליים, אך כיוונון עדין (\textenglish{fine-tuning}) נדרש להתאמה למשימה ספציפית. בכיוונון עדין מפוקח (\textenglish{supervised fine-tuning}), מתווסף ראש סיווג (\textenglish{classification head}) ומותאמים כלל פרמטרי המודל על גבי נתוני המשימה. גישות כיוונון עדין פרמטרי-יעיל (\textenglish{PEFT}) כגון \textenglish{LoRA} מוסיפות עדכון בדרגה נמוכה:

\begin{equation}
\mathbf{W}' = \mathbf{W} + \mathbf{B}\mathbf{A}, \qquad \mathbf{B} \in \mathbb{R}^{d \times r},\; \mathbf{A} \in \mathbb{R}^{r \times k},\; r \ll \min(d,k)
\end{equation}

גישה זו מצמצמת את מספר הפרמטרים הניתנים לאימון תוך שמירה על ביצועים תחרותיים. גישת \textenglish{prompt-based learning} מנצלת את יכולות ה-\textenglish{in-context learning} ללא עדכון פרמטרים כלל. העברת למידה חוצת-שפות עם \textenglish{mBERT} \cite{pires2019multilingual} הדגימה כי מודל שעבר כיוונון על אנגלית מסוגל לבצע משימות בעברית ובשפות אחרות ללא נתוני כיוונון ייעודיים.

\section{מודלי שפה גדולים ויכולות מתפתחות}

חוקי ה-\textenglish{scaling} \cite{brown2020language} מראים כי ביצועי המודל מתפתחים בצורה חזויה עם כמות הפרמטרים \(N\), כמות הנתונים \(D\) ועוצמת המחשוב \(C\):

\begin{equation}
L(N) \propto N^{-\alpha}, \qquad L(D) \propto D^{-\beta}, \qquad L(C) \propto C^{-\gamma}
\end{equation}

מעבר לשיפורים הצפויים, \textenglish{scaling} מוביל לתופעת \textenglish{emergent capabilities} --- יכולות המופיעות בצורה פתאומית מעל סף גודל מסוים, כגון חשבון רב-שלבי, מעקב הוראות מורכב, וכתיבת קוד. מתודולוגיית \textenglish{RLHF} (\textenglish{Reinforcement Learning from Human Feedback}) מיישרת מודלים עם כוונות המשתמש:

\begin{equation}
\mathcal{J}(\theta) = \mathbb{E}\!\left[r_\phi(\mathbf{x}, \mathbf{y}) - \beta \, \mathrm{KL}\!\left(\pi_\theta(\mathbf{y}\mid\mathbf{x}) \,\|\, \pi_{\text{ref}}(\mathbf{y}\mid\mathbf{x})\right)\right]
\end{equation}

כאשר \(r_\phi\) הוא מודל התגמול, \(\pi_\theta\) היא מדיניות המודל המכוונן, ו-\(\pi_{\text{ref}}\) היא מדיניות הייחוס. גישה זו הובילה לפיתוח מודלים מיישורים כגון \textenglish{InstructGPT} ו-\textenglish{ChatGPT}, המשלבים יכולות גנרטיביות גבוהות עם התאמה לכוונות המשתמש. בסיכום, שדה מודלי השפה הגדולים נמצא בהתפתחות מואצת, ומציב תשתית חיונית להבנת יישומי הטרנספורמר לעיבוד שפה עברית שיידונו בפרק הבא.
